\chapter*{谢辞}
\addcontentsline{toc}{chapter}{谢辞}

行文至此，毕业设计将近尾声，四年的本科生活也即将落笔成章。谨以此篇，献给在我求学路上给予过帮助、支持与善意的每一位老师、同学、朋友和家人，也献给这段从天山脚下到同济的青春岁月。

首先，衷心感谢指导老师徐凡老师在毕业设计全过程中的悉心指导。从课题选定、研究思路梳理，到实验方案设计、结果分析和论文结构修改，徐凡老师都给予了我耐心而具体的建议。尤其是在论文多轮修改过程中，老师对摘要表述、章节逻辑、图表格式和结论边界提出了细致意见，使我更加认识到，工程研究不仅需要得到实验结果，也需要以严谨、克制和清晰的方式表达结果。老师严谨认真的治学态度，使我受益良多；这份训练也不会止于论文完成之时，而会继续影响我今后的学习、工作与人生选择。

感谢同济大学电子与信息工程学院在本科阶段为我提供的学习环境、课程训练和科研资源。作为微电子科学与工程专业的学生，四年的学习经历使我逐步建立起工程建模、程序实现、实验分析和论文写作的基本能力，也为本次毕业设计的完成奠定了基础。从四平南校区到嘉定校园，从课堂学习到实验实践，同济大学严谨求实的学风，让我明白真正的成长并不来自一时的顺遂，而来自长期、踏实、反复的积累。嘉定图书馆里安静的自习时光，嘉定友园中平凡而充实的生活片段，以及嘉定校园四季更替中的奔走与停留，都将成为我大学记忆中难以忘却的一部分。

回望来路，我来自新疆昌吉，来自天山脚下的一座小城。年少时看惯了辽阔的天空、远处的雪山和故乡的风，后来一路向东，来到上海求学。地域的距离，不只是地图上漫长的迁徙，也意味着视野、生活方式和内心秩序的重新建立。初到上海时，我曾对陌生的城市、紧凑的节奏和复杂的专业课程感到不适应；在后来的学习与毕业设计中，也曾因为实验结果不理想、代码调试反复、论文修改多次而陷入焦虑。但这些经历最终让我明白，成长并不是把所有路都走得漂亮，而是在迷茫和受挫之后，仍愿意整理自己，继续向前。

拿破仑曾说：“生活的成果不在于永不失败，而在于屡仆屡起。”这句话曾在许多低落的时刻提醒我：失败本身并不可怕，重要的是是否还有重新站起、重新开始的勇气。如今再看这段本科求学经历，那些曾经让我困惑、怀疑甚至疲惫的时刻，也都成为塑造我的一部分。它们让我更加明白严谨的意义，也让我学会在不确定中保持耐心，在平凡的努力中积蓄力量。

感谢大学期间以及此前相识的朋友们。感谢你们在学习、生活和情绪低落时给予我的陪伴与鼓励。许多普通的日子，因为你们的理解和支持而变得温暖；许多看似难以跨过的阶段，也因为有人同行而不再孤单。青春里有太多无法一一写尽的人和事，所念诸多，言有不尽，唯愿我们在各自奔赴的道路上，都能保有真诚、热烈与笃定。

最后，感谢我的父母和家人长期以来给予我的理解、支持和鼓励。从新疆昌吉到上海同济，求学路途遥远，而你们始终是我身后最稳定的力量。正是你们无声而持续的支持，使我能够专注完成本科阶段的学习和毕业设计工作。也感谢学校的工作人员，以及在生活中曾经帮助过我的每一个人。你们或许只是给予过一句提醒、一次帮助、一个微小的善意，却都在某些具体的时刻支撑我继续前行。

谨以此文，向所有关心、帮助和支持过我的人表示衷心感谢。未来无论身处何处，我都会铭记这段从天山脚下到同济的求学经历，带着感恩、严谨与屡仆屡起的心态，继续走向人生的下一程。
